% ------------------------------------------------------------------ cover
\thispagestyle{blank}
\begin{tikzpicture}[remember picture,overlay]
  \fill[disc] (current page.south west) rectangle (current page.north east);
  % the 8088's twenty address lines, A19 to A0
  \addressbandlabels{(current page.north west)+(0,{-3.05in-\bleed})}{\paperwidth}{0.42in}{white!70!disc}
  \node[anchor=north,text=white!70!disc,inner sep=0pt,font=\sffamily\small] at ($(current page.north)+(0,{-3.47in-\bleed})$) {the 8088's twenty address lines: one megabyte};
  % The Lyon Mountain Company's lion, alone, small, upper right
  \node[anchor=north east,inner sep=0pt] at ($(current page.north east)+({-0.8in-\bleed},{-0.55in-\bleed})$)
    {\includegraphics[height=1.0in]{figures/lmc-lion.png}};
  \node[anchor=north west,text=white,inner sep=0pt,align=left] at ($(current page.north west)+({0.85in+\bleed},{-3.95in-\bleed})$)
    {\cond\bfseries\fontsize{68}{68}\selectfont RetrOS/LT};
  \node[anchor=north west,text=white,inner sep=0pt,align=left,font=\cond\fontsize{18}{22}\selectfont] at ($(current page.north west)+({0.9in+\bleed},{-5.0in-\bleed})$)
    {\subtitlelines};
  \node[anchor=north west,text=white,inner sep=0pt,align=left] at ($(current page.north west)+({0.9in+\bleed},{-5.95in-\bleed})$)
    {\sffamily\bfseries\fontsize{24}{28}\selectfont \booktitle};
  \node[anchor=north west,text=white!80!disc,inner sep=0pt,align=left,font=\sffamily\fontsize{11.5}{15}\selectfont] at ($(current page.north west)+({0.9in+\bleed},{-6.5in-\bleed})$)
    {RetrOS/LT Version 1.0, DOS 2.10\\for the IBM Personal Computer, the XT, and compatibles};
  % the processor
  \node[anchor=south,inner sep=0pt] at ($(current page.south)+(0,{1.25in+\bleed})$)
    {\dipchip{3.1}};
\end{tikzpicture}
\newpage
% Two-sided: the inside of the front cover is left empty.
\ifmwtwoside\thispagestyle{blank}\null\newpage\fi

% ------------------------------------------------------------------ title page
\frontpages
\setcounter{page}{1}
\pagestyle{blank}
\vspace*{0.4in}
\hspace*{-\dimexpr\marginparwidth+\marginparsep\relax}%
\begin{minipage}[t]{\dimexpr\textwidth+\marginparwidth+\marginparsep\relax}
\raggedright
{\cond\bfseries\fontsize{40}{44}\selectfont\color{disc} RetrOS/LT\par}
\vspace{6pt}
{\cond\fontsize{17}{21}\selectfont\color{disc} \subtitleline\par}
\vspace{14pt}
\begin{tikzpicture}\addressband{(0,0)}{\linewidth}{4pt}\end{tikzpicture}\par
\vspace{14pt}
{\sffamily\bfseries\fontsize{22}{26}\selectfont \booktitle\par}
\vspace{8pt}
{\large RetrOS/LT Version 1.0, for the IBM Personal Computer,\\
the Personal Computer XT, and compatible computers\par}
\vspace{0.5in}
\hspace*{0.3in}\dipchip{2.0}\par
\vspace{24pt}
{\small\textbf{The Lyon Mountain Company}\\
Republic Lyon Mountain Corporation\\
Lyon Mountain, New York 12952 USA\par}
\end{minipage}
\newpage

% ------------------------------------------------------------------ edition notice
\pagestyle{blank}
\vspace*{\fill}
\begin{wide}
\small
\textbf{\edition}

This edition applies to Version 1.0 of RetrOS/LT, build \release, and to
all later releases until otherwise indicated in new editions. Changes are
made periodically to the information herein. \editionchanges

A form for readers' comments is at the back of this book. If the form has
been removed, comments may be addressed to P.O. Box 290, Lyon Mountain,
New York 12952 USA.

\textbf{Where RetrOS/LT comes from.} RetrOS/LT is derived from the MS-DOS 2
source code that Microsoft Corporation released under the MIT License.
That code is Copyright \copyright\ Microsoft Corporation; the notice that
must accompany it is printed in Appendix~D and is on the RetrOS/LT
diskettes. RetrOS/LT is not MS-DOS and is not PC DOS. It is an independent
product of The Lyon Mountain Company, and is not affiliated with, endorsed
by, or supported by Microsoft Corporation or International Business
Machines Corporation.

\textbf{RetrOS/LT is licensed, not sold}, under the RetrOS/LT License
printed in Appendix~D: a Single-Machine License, for one computer at a
time, or a Blanket License, for all the computers one owner owns or
leases.

\textbf{Trademarks.} IBM, Personal Computer XT and PC DOS are trademarks
of International Business Machines Corporation. Microsoft and MS-DOS are
trademarks of Microsoft Corporation. Intel is a trademark of Intel
Corporation. Other names used in this book may be trademarks of their
owners.

\copyright\ Copyright The Lyon Mountain Company, Republic Lyon Mountain Corporation 2026
\end{wide}
\newpage

% ------------------------------------------------------------------ preface
\pagestyle{plain}
\renewcommand\runningtitle{About This Book}
\chapter*{About This Book}

\addcontentsline{toc}{chapter}{About This Book}

This book tells you how to install and use RetrOS/LT, a disk operating
system for the IBM Personal Computer, the Personal Computer XT, and
compatible computers. RetrOS/LT is the small member of the RetrOS family:
where RetrOS/86 and RetrOS/286 carry DOS 4 to the PC and AT, RetrOS/LT is
DOS 2.10, as the IBM Personal Computer had it, made again from Microsoft's
source.

The address band at the head of each Section is the 8088's twenty address
lines, A19 to A0: the one megabyte RetrOS/LT and every program it runs
live in.

\section{How This Book Is Organized}

Read Section~1 first. If you are installing RetrOS/LT on a fixed disk,
read Section~2. If you have used PC DOS 2, you already know RetrOS/LT; go
straight to Section~6, the command reference, and read the rest when you
need it.

\begin{tabularx}{\textwidth}{@{}l>{\raggedright\arraybackslash}X@{}}
\toprule
\textbf{Section} & \textbf{Read it if you}\\
\midrule
1 Introduction & are new to RetrOS/LT, or want to know what it runs on\\
2 Installing RetrOS/LT & are putting RetrOS/LT on a fixed disk\\
3 Using RetrOS/LT & are new to DOS: the prompt, files, directories and running programs\\
4 Working with Disks & format, check, copy or back up diskettes and fixed disks\\
5 Setting Up Your System & want to change \file{CONFIG.SYS} or \file{AUTOEXEC.BAT}, load a device driver, or set up a printer or display\\
6 Commands & want the format and switches of a command\\
7 Problem Solving & have a problem, or see a message you do not understand\\
Appendixes & want the list of files, the technical details, or the license\\
\bottomrule
\end{tabularx}

A Glossary and an Index follow the Appendixes. At the very back are a Quick
Reference Card and a Reader's Comment Form.

\section{Page Numbers}

Pages are numbered within each Section: page 4-7 is the seventh page of
Section~4, and page B-2 the second page of Appendix~B. Tables are numbered
the same way.

\section{Conventions}

Keys you press are drawn as keys: \key{Enter}, \key{Y}, \key{N}. When keys
are pressed together they are joined: \key{Ctrl}\,+\,\key{Alt}\,+\,\key{Del}
means hold down \key{Ctrl} and \key{Alt}, and press \key{Del}. The key
the IBM keyboard marks with a bent arrow is called \key{Enter} here.

What you type is printed \typed{like this}. What the computer shows is
printed \shown{like this}:

\begin{dialogue}
A>\textbf{VER}

RetrOS/LT Version 1.0, DOS 2.10
\end{dialogue}

In the format of a command, words in capitals are typed as shown (in
capitals or small letters, as you like); words in \var{italics} stand for
something you supply; items in square brackets [\ ] may be left out; and a
bar | separates choices of which you type one. You do not type the
brackets or the bar.

\note{A note points out something useful, or something easy to overlook.}

\caution{A caution warns you of something that could lose your files or
keep your computer from starting.}

\mwnewpart
\renewcommand\runningtitle{Contents}
\renewcommand\contentsname{Contents}
\tableofcontents
